%======================================================================
\section{Introduction}\label{sec:intro}
%======================================================================

Mantles and other planetary interiors are heated from within by radioactive decay. Their viscosity is so
large that the Prandtl number is effectively infinite (about $10^{23}$ for the Earth's mantle). The idealised
model is uniform internally heated convection between two isothermal plates. Without flow, half the heat
leaves through each plate. Convection breaks this symmetry: hot fluid rises, plumes of cold fluid descend from
the unstable upper boundary layer, and a stably stratified layer forms at the base. The fraction $F_B$ of the
heat that still leaves through the base measures how strongly convection redistributes heat upward. It plays a
role similar to that of the Nusselt number in Rayleigh--B\'enard convection \citep{Goluskin2016}.

Upper bound theory asks how small $F_B$ can be for \emph{any} solution, steady or turbulent, as the flux
Rayleigh number $R$ grows. That $F_B\ge0$ follows from the minimum principle \citep{GoluskinSpiegel2012}.
Quantitative bounds came much later. At arbitrary Prandtl number the best results are exponentially small in
$R$ \citep{ArslanEtAl2021,KumarEtAl2022}. At infinite Prandtl number, \citet{ArslanEtAl2023} proved
$F_B\gtrsim R^{-2}$, and \citet{ArslanRojas2025} improved this to
\[
F_B\ \gtrsim\ R^{-2/3}\quad\text{(no-slip)},\qquad F_B\ \gtrsim\ R^{-40/29}\quad\text{(stress-free)},
\]
using the auxiliary-functional (background) method with the minimum principle imposed as a constraint. The
companion bound on the mean temperature at infinite Prandtl number is $\langle T\rangle\ge0.419(R\log R)^{-1/4}$
\citep{WhiteheadDoering2011}, and rotating layers have also been treated \citep{Arslan2024rot,ZhangEtAl2026}.
\citet{ChenEtAl2024} obtained the faster rate $F_B\ge 8.842R^{-1/3}$ at any Prandtl number, but only
\emph{conditionally}. Their Eq.~(22), $\langle wT\rangle\le(\frac14-8.842R^{-1/3})^{1/2}$, combines the rigorous
bound $\langle T\rangle\ge4.421R^{-1/3}$ of \citet{LuDoeringBusse2004} with the inequality
$\langle T\rangle\le\frac18-\frac12\langle wT\rangle^2$ (their Eq.~21). That inequality rests on the hypothesis
(their Eq.~14) that the mean temperature profile lies below the triangle $\min\{F_Bz,(1-z)(1-F_B)\}$, which they
justify by a Taylor expansion at the walls and by simulations, and which they list among the limitations of their
work. (Their text gives the approach of $\langle wT\rangle$ to $\tfrac12$ as $R^{-1/6}$; since
$\tfrac12-\sqrt{1/4-x}\ge x$, Eq.~(22) in fact implies $F_B\ge8.842R^{-1/3}$.) The two-dimensional
simulations of \citet{ArslanRojas2025} give $F_B\sim R^{-0.092}$ and $R^{-0.12}$. They observe that ``whether any flow
saturates the known bounds on $F_B$ remains an open question'' and that further optimisation could improve the
bounds. At finite Prandtl number, \citet{GoluskinvanderPoel2016} found that the decrease of $F_B$ with $R$
reverses in the turbulent regime in two dimensions but continues in three. The problem has also been posed as item~3 of a recent list of open problems in geology suited to
frontier models. That list notes, by contrast, that the basally heated (Rayleigh--B\'enard) counterpart is now
sharp. At infinite Prandtl number the bounds $\mathrm{Nu}\lesssim\mathrm{Ra}^{1/3}(\ln\mathrm{Ra})^{2/3}$
\citep{ConstantinDoering1999}, $\mathrm{Ra}^{2/5}$ from the background method \citep{DoeringConstantin2001},
$\mathrm{Ra}^{1/3}(\ln\mathrm{Ra})^{1/3}$ \citep{DoeringOttoReznikoff2006} and
$\mathrm{Ra}^{1/3}(\ln\ln\mathrm{Ra})^{1/3}$ \citep{OttoSeis2011} were followed by $\mathrm{Nu}\lesssim\mathrm{Ra}^{1/3}$
without logarithms \citep{ChanilloMalchiodi2025}, extended to Navier-slip walls by \citet{Seis2026}. It asks for two things: sharper bounds in the internally heated case, or
flows that attain them; and a constructed solution of the basally heated problem with
$\mathrm{Nu}\ge c\,\mathrm{Ra}^{1/3}$.

\paragraph{Results.} Our main contribution is to the first question. We prove
(Theorem~\ref{thm:main} and Corollary~\ref{cor:R}) that every classical solution satisfies
\[
\overline{F_B}\ \ge\ c\,(R\mathcal T\ln R\mathcal T)^{-1/3}\quad\text{(no-slip)},\qquad
\overline{F_B}\ \ge\ c\,(R\mathcal T)^{-1/2}\quad\text{(stress-free)},
\]
where $\mathcal T$ is the long-time maximum temperature, and that $\mathcal T\le C(1+\ln|\Omega|)$ for every
incompressible flow (Theorem~\ref{thm:Linfty}). At fixed aspect ratio this improves the exponents of the best unconditional bounds from
$2/3$ to $1/3$ (up to a logarithm) and from $40/29$ to $1/2$. In particular, at infinite Prandtl number the
exponent $1/3$ of the conditional bound of \citet{ChenEtAl2024} is obtained, up to the logarithm, without their
profile hypothesis. Our computed flows also let us test that hypothesis directly (Section~\ref{sec:numerics}). The
flux argument does not use the energy identities on which the background method rests; only the temperature bound
uses standard $L^p$ estimates. It has three steps:
\begin{enumerate}
\item Because $T\ge0$, the horizontally averaged convective flux at height $z$ is at most the largest vertical
speed at that height times the mean temperature there. This inequality is \emph{linear} in the temperature.
\item Near the base the infinite-Prandtl velocity obeys $|w|\le\tfrac12R\mathcal T\,\Phi(z)$. The wall
function $\Phi$ is an explicit integral of $L^1$ norms of Green's-function kernels, with
$\Phi(z)=\tfrac12\nu_2z^2\ln(1/z)+O(z^2)$ for no-slip and $\Phi(z)\sim c_{\rm w}z$ for stress-free walls.
\item Steps 1--2 turn the horizontally averaged heat balance into the differential inequality
$\Theta'\le F_B-z+m(z)\Theta$ for the mean temperature profile $\Theta\ge0$, whose integrating factor gives an
explicit lower bound on $F_B$.
\end{enumerate}
We also determine the exact large-$R$ asymptotics of the bound functional, which give the best constants that
inequality \eqref{eq:mainbound} can yield (Proposition~\ref{prop:sharp}). The constants depend on two norms of a
slab correction to the half-space Green's function. We prove that they are finite and evaluate them numerically,
with numerical sensitivity checks, not uniform error enclosures and a second, more conservative route, but without interval certification
(Section~\ref{sec:status} lists what is proved and what is computed).

For the second half of each question (attainment and construction) we offer computations, not theorems. We
solve the full infinite-Prandtl equations in two dimensions with a Fourier--Chebyshev time-stepper and a Newton
solver for steady rolls, both validated against published linear thresholds and benchmark values to within
the precision quoted for them.
\begin{itemize}
\item In simulated flows we test each inequality of the proof directly. Both hold with large margins, and the
measured fluxes stay far above the bound. Near the base the convective flux is close to zero and in places
negative, the reverse of what the extremal configuration needs.
\item Steady internally heated rolls with the period chosen to minimise $F_B$ do essentially no better than the
time-dependent simulations of \citet{ArslanRojas2025}: $F_B\approx1.29R^{-0.093}$ for
$1.6\times10^6\le R\le2.8\times10^7$ (Section~\ref{sec:steady}).
\item For Rayleigh--B\'enard convection at infinite Prandtl number, within the family of mirror-symmetric,
layer-spanning rolls and the period range we scanned, the wide-roll maximum of Nu over the period, which yields
$\mathrm{Ra}^{1/3}$ at $\mathrm{Pr}=1$ \citep{WenGoluskinDoering2022}, is not detected above the sampled interval between $\mathrm{Ra}=10^5$ and
$1.8\times10^5$. Above that, the largest maxima found are narrow rolls, which follow
$\mathrm{Nu}\approx0.153\,\mathrm{Ra}^{0.298}$ with period $\Gamma^*\approx11.7\,\mathrm{Ra}^{-1/4}$ up to
$\mathrm{Ra}=10^8$ (Section~\ref{sec:steady}). This extends to infinite Prandtl number a trend that
\citet{SondakSmithWaleffe2015} found at $\mathrm{Pr}=10$ and $100$, where the narrow-roll maximum eventually
overtakes the wide-roll one, and it matches the $\mathrm{Ra}^{3/10}$ scaling that \citet{WenEtAl2026} derived for
high-wavenumber rolls between stress-free walls. If the trend persists, two-dimensional steady rolls of this
family will not attain the upper bound at infinite Prandtl number.
\end{itemize}
A rigorous construction remains open, and we explain why the structure of these rolls makes it hard.

\paragraph{Relation to earlier work.} The pointwise route to bounds at infinite Prandtl number goes back to
\citet{ConstantinDoering1999}. They combined the maximum principle $0\le T\le1$ with a logarithmic estimate for
the Stokes kernel to bound the Nusselt number in Rayleigh--B\'enard convection. For internally heated
convection every published bound on $F_B$ uses the background method. Even the recent work that uses
pointwise velocity estimates \citep{Arslan2024rot} applies them to individual Fourier modes inside a quadratic
spectral constraint. The integrating-factor inequality (Lemma~\ref{lem:IF}) and the resulting exponents were
not found in the literature search recorded in the accompanying novelty report. That search traced the citation
graphs of the internally heated bound papers and ran full-text and arXiv listing searches. Our temperature
bound is a time-dependent slab variant of known $L^\infty$ estimates for incompressible drifts
\citep{BerestyckiEtAl2010,IyerEtAl2010}.
